\section{Spectral applications}\label{sec:spectral}

Throughout this section let
\[
 T=a_0+\sum_{j=1}^n a_j\otimes s_j\in M_m(\M),\qquad
 S_z=T-z1_m,\qquad a_z=a_0-z1_m,
\]
with arbitrary complex coefficient matrices.
We keep the maps $\eta$ and $\rho$ from \eqref{eq:coefficient-maps}.
Applying Theorem~\ref{thm:main} to $S_z$ gives the logarithmic potential of $\mu_T$ and characterizes $\spec(T)$.

\subsection{Shifted potentials and spectral interpretation}
For $b>0$, write $Z_b(z)=\mathcal T_b^{-1}(a_z)$ and
\[
 \mathcal V_z(b)=\log\det H_b-\log\det b-m
                   +\operatorname{Re}\Tr(a_z^*Z_b(z)),
 \qquad H_b=1_m+\eta(b).
\]
Theorem~\ref{thm:main} gives
\begin{equation}\label{eq:shifted-potential}
 u(z):=\log\Delta(S_z)=\frac1{2m}\inf_{b>0}\mathcal V_z(b),
\end{equation}
with the value $-\infty$ permitted.
It also gives
\begin{equation}\label{eq:full-spectrum}
 z\notin\spec(T)
 \quad\Longleftrightarrow\quad
 \mathcal V_z\text{ has a critical point}
 \quad\Longleftrightarrow\quad
 \inf_{b>0}\mathcal V_z(b)\text{ is attained}.
\end{equation}
At a point of the resolvent set, the unique minimizer is $b_z=E_m((S_z^*S_z)^{-1})$.
The criterion concerns bounded invertibility; finiteness of $u(z)$ alone does not imply it.

We use the Wirtinger derivatives $\partial_z=(\partial_x-i\partial_y)/2$ and $\partial_{\bar z}=(\partial_x+i\partial_y)/2$, where $z=x+iy$.
\begin{definition}[Brown measure]\label{def:brown}\label{def:brown-measure}
The Brown measure of $T$ with respect to $\varphi_m$ is the probability measure
\begin{equation}\label{eq:brown-definition}
 \mu_T=\frac1{2\pi}\nabla_z^2u,\qquad
 \nabla_z^2=4\partial_{\bar z}\partial_z,
\end{equation}
where the Laplacian is taken in the sense of distributions.
\end{definition}
Here $u$ is subharmonic and locally integrable, and $\supp\mu_T\subseteq\spec(T)$.
With the normalization at infinity fixed by the determinant, it is exactly the logarithmic potential
\begin{equation}\label{eq:brown-logarithmic-potential}
 u(z)=\int_{\C}\log|z-w|\dd\mu_T(w).
\end{equation}
These properties of the Brown measure are recalled in \cite[Chapter 11, Theorem 5, Definition 6 and (11.6)--(11.8)]{MS17}.
We use them to interpret the variational formula; no equality between $\supp\mu_T$ and $\spec(T)$ is assumed.

For $\varepsilon>0$, set
\begin{equation}\label{eq:spectral-potential}
 u_\varepsilon(z)=\tfrac12\varphi_m
                   \log(S_z^*S_z+\varepsilon^21_m).
\end{equation}
The rescaled minimum \eqref{eq:scaled-minimum} reads
\begin{equation}\label{eq:brown-regularized-minimum}
 u_\varepsilon(z)=\frac1{2m}\min_{b>0}F_\varepsilon(z,b),
 \qquad F_\varepsilon(z,b)=\mathcal V_z(b)+\varepsilon^2\Tr b.
\end{equation}
Its unique minimizer $b_\varepsilon(z)$ and the matrix $Z_\varepsilon(z)=\mathcal T_{b_\varepsilon(z)}^{-1}(a_z)$ are
\begin{equation}\label{eq:brown-regularized-optimizer}
 \begin{split}
 b_\varepsilon(z)&=E_m((S_z^*S_z+\varepsilon^21_m)^{-1}),\\
 Z_\varepsilon(z)&=Z_{b_\varepsilon(z)}(z)
     =E_m((S_zS_z^*+\varepsilon^21_m)^{-1}S_z).
 \end{split}
\end{equation}
We will use these matrices both to recover radial masses and to compute a regularized density.
The regularized Cauchy transform and its relation to hermitization are developed in \cite[arXiv version, Sections~2.1.6 and~2.2]{BSS18}.
Here \eqref{eq:brown-regularized-minimum} expresses the same regularized potential through a positive-matrix minimum.

\subsection{Determinant lower bounds and atoms of the Brown measure}
\label{subsec:brown-atoms}
We first prove Corollary~\ref{cor:brown-capacity}, then locate the possible atoms without assuming positive capacity.
For $a\in M_m(\C)$ and $b>0$, put $q_a(b)=\operatorname{Re}\Tr(a^*\mathcal T_b^{-1}(a))$.

% Proof inputs: Yin_C_F0126 and Yin_C_F0172.
\begin{proof}[Proof of Corollary~\ref{cor:brown-capacity}]
Use the real inner product $\langle X,Y\rangle_{\R}=\operatorname{Re}\Tr(X^*Y)$.
The positivity established in \eqref{eq:quadratic-lower-bound}, at $\varepsilon=1$, implies
\[
 q_a(b)=\langle a,\mathcal T_b^{-1}(a)\rangle_{\R}\ge0.
\]
Put $T_0=\sum_j a_j\otimes s_j$.
Taking infima in Theorem~\ref{thm:main} gives, for every $a\in M_m(\C)$,
\begin{equation}\label{eq:brown-capacity-lower-bound}
 \Delta(a+T_0)\ge\Delta(T_0)
               =e^{-1/2}\Cap(\eta)^{1/(2m)}.
\end{equation}
Taking $a=a_0-z1_m$ proves \eqref{eq:intro-uniform-bound}.
At positive capacity, this gives the uniform bound
\[
 u(z)\ge -\frac12+\frac1{2m}\log\Cap(\eta)>-\infty.
\]
An atom at $z$ would make \eqref{eq:brown-logarithmic-potential} equal $-\infty$ there, since its positive part is bounded on the compact support.
This proves the assertion about atoms.
\end{proof}

% Proof input: Yin_L_F0042, Section 6.
\begin{proposition}[Possible atom locations]\label{prop:brown-atoms}
For every $z\in\C$,
\begin{equation}\label{eq:brown-atom-locations}
 \det(a_0-z1_m)\ne0\ \Longrightarrow\ \Delta(T-z1_m)>0.
\end{equation}
In particular every atom of $\mu_T$ belongs to $\spec(a_0)$.
\end{proposition}
\begin{proof}
Put $T_0=\sum_j a_j\otimes s_j$, and let $u_0$ be the logarithmic potential of $\mu_{T_0}$.
Applying Theorem~\ref{thm:main} to $T_0-re^{i\theta}1_m$ and using the real-linearity of $\mathcal T_b^{-1}$, we obtain
\begin{equation}\label{eq:brown-ray-monotonicity}
 u_0(re^{i\theta})=\frac1{2m}\inf_{b>0}
 \{A(b)+r^2q_{-e^{i\theta}1_m}(b)\},\qquad
 A(b)=\log\det H_b-\log\det b-m.
\end{equation}
Since $q_{-e^{i\theta}1_m}(b)\ge0$ by \eqref{eq:quadratic-lower-bound}, the potential is nondecreasing in $r\ge0$ for each fixed $\theta$.
It cannot take the value $-\infty$ at a positive radius.
Indeed, for any bounded real interval $I$ and any $w\in\C$,
\[
 \int_I\log^-|re^{i\theta}-w|\dd r
 \le\int_{-1}^1-\log|t|\dd t=2,
 \qquad \log^-t=\max\{-\log t,0\}.
\]
The inequality follows from $|re^{i\theta}-w|\ge|r-\operatorname{Re}(e^{-i\theta}w)|$.
The positive logarithmic part is uniformly bounded when $w$ ranges over a compact set.
Fubini's theorem therefore shows that a compactly supported logarithmic potential is integrable on every finite line segment.
If $u_0(r_0e^{i\theta})=-\infty$ for $r_0>0$, monotonicity would make it $-\infty$ on $[r_0/2,r_0]e^{i\theta}$, a contradiction.
Hence $\Delta(T_0-z1_m)>0$ for $z\ne0$.

Now suppose $A=a_0-z1_m$ is invertible and write
\[
 T-z1_m=A(1_m+C),\qquad
 C=\sum_j A^{-1}a_j\otimes s_j.
\]
The centered conclusion applied to $C$ at the shift $-1$ gives $\log\Delta(1_m+C)>-\infty$.
If $c>0$ is the smallest singular value of $A$, then, with $B=1_m+C$, $(AB)^*(AB)\ge c^2B^*B$.
Add $\varepsilon^21_m$, apply the operator-monotone logarithm and the trace, and let $\varepsilon\downarrow0$.
The result is
\[
 \log\Delta(T-z1_m)\ge\log c+\log\Delta(1_m+C)>-\infty.
\]
This proves \eqref{eq:brown-atom-locations} and the last assertion.
\end{proof}

The atom locations in Proposition~\ref{prop:brown-atoms} need not all carry positive mass.
Nor does atomlessness imply absolute continuity with respect to planar area.
For example, if $T=s$ is a standard selfadjoint semicircular element, then $\eta=\id$ and $\Cap(\eta)=1$, while $\mu_T$ is supported on $\R$.

A second distinction concerns regularization at a fixed point.
Spectral calculus gives
\begin{equation}\label{eq:brown-kernel-limit}
 \lim_{\varepsilon\downarrow0}\varepsilon^2
               \tr_m b_\varepsilon(z)
   =\varphi_m(1_{\{0\}}(S_z^*S_z)).
\end{equation}
This is the trace of the kernel projection of $S_z$, which need not equal $\mu_T(\{z\})$.
More generally, the former is at most the latter; see \cite[Section~2.1, p.~12]{ACSY24} for this distinction between point-spectrum weights and atoms of the Brown measure.
For the deterministic matrix $T=\left(\begin{smallmatrix}0&1\\0&0\end{smallmatrix}\right)$, the right side at zero is $1/2$, whereas $\mu_T=\delta_0$: indeed $\log\Delta(T-z1_2)=\log|z|$.
Thus divergence of $b_\varepsilon(z)$ alone does not determine $\mu_T(\{z\})$.

\subsection{Centered pencils with vanishing \texorpdfstring{$\rho$}{rho}}
\label{subsec:brown-radial}
Assume in this subsection that $a_0=0$ and $\rho=0$ as a map on $M_m(\C)$.

\begin{example}[Matrix coefficients and free circular elements]
\label{ex:brown-circular-coefficients}
Let $s_1,\ldots,s_{2r}$ be free standard semicircular elements and set
\[
 c_k=\frac{s_{2k-1}+is_{2k}}{\sqrt2},\qquad 1\le k\le r.
\]
Then $c_1,\ldots,c_r$ are free standard circular elements.
For matrices $d_1,\ldots,d_r\in M_m(\C)$, the operator
\[
 T=\sum_{k=1}^r d_k\otimes c_k
\]
is a centered semicircular pencil with coefficients
\[
 a_{2k-1}=\frac{d_k}{\sqrt2},\qquad
 a_{2k}=\frac{id_k}{\sqrt2}.
\]
For every $x\in M_m(\C)$, these pairs give
\[
 \rho(x)=\frac12\sum_{k=1}^r
           \bigl(d_kxd_k+(id_k)x(id_k)\bigr)=0,
 \qquad
 \eta(x)=\sum_{k=1}^r d_kxd_k^*.
\]
Thus Proposition~\ref{prop:brown-radial} applies to $T$ with this map $\eta$.
\end{example}

For a centered semicircular pencil, all $M_m(\C)$-valued cumulants of $(T,T^*)$ except those of order two vanish.
The condition $\rho=0$ also gives $E_m(TxT)=E_m(T^*xT^*)=0$.
Its only possibly nonzero cumulants are therefore
\[
 E_m(TxT^*)=\eta(x),\qquad E_m(T^*xT)=\eta^*(x),
 \qquad x\in M_m(\C).
\]
By \cite[Theorem 2.9(b)]{DG24}, $T$ is $M_m(\C)$-valued $R$-diagonal.
Since all cumulants of order greater than two vanish, \cite[Definition 2.10]{DG24} identifies $T$ as an $M_m(\C)$-valued circular element.
These cumulants are unchanged when $T$ is replaced by $e^{i\theta}T$, so its $M_m(\C)$-valued $*$-distribution is rotation invariant.
We now apply Theorem~\ref{thm:main} to $T-z1_m$ to prove Proposition~\ref{prop:brown-radial} and the following formulas for the logarithmic potential and density of $\mu_T$.

% Proof input: fixed radial deduction Yin_C_F0170.
\begin{proposition}[Radial density and minimizers]\label{prop:brown-radial-details}
The function $f$ in Proposition~\ref{prop:brown-radial} is given by
\begin{equation}\label{eq:brown-radial-envelope}
 f(s)=\frac1{2m}\inf_{b>0}
 \left\{\log\det H_b-\log\det b-m
                       +s\Tr(H_b^{-1}b)\right\}.
\end{equation}
The radial distribution $F$ in \eqref{eq:brown-radial-cdf} is locally Lipschitz on $(0,\infty)$, and
\begin{equation}\label{eq:brown-radial-density}
 \mu_T=\alpha\delta_0+\beta(z)\dd^2z,\qquad
 \beta(z)=\frac1\pi F'(|z|^2),\qquad
 0\le\beta(z)\le\frac{F(|z|^2)}{\pi|z|^2}
              \le\frac1{\pi|z|^2}\quad\text{a.e. }z\ne0,
\end{equation}
where $\alpha=\lim_{s\downarrow0}F(s)$.
For every $s>0$,
\begin{equation}\label{eq:brown-radial-optimizer}
 F(s)=\frac{s}{m}\lim_{\varepsilon\downarrow0}
 \Tr\bigl((1_m+\eta(b_\varepsilon(s)))^{-1}b_\varepsilon(s)\bigr),
\end{equation}
where the minimizer in \eqref{eq:brown-regularized-optimizer} depends only on $s=|z|^2$.
\end{proposition}
\begin{proof}[Proof of Propositions~\ref{prop:brown-radial} and
\ref{prop:brown-radial-details}]
Since $\rho=0$, the solution of $\mathcal T_b(Z_b(z))=-z1_m$ is $Z_b(z)=-zH_b^{-1}b$.
Substitution in \eqref{eq:shifted-potential} proves \eqref{eq:brown-radial-envelope}.
Its slope at each fixed trial is positive, so $f$ is nondecreasing and concave.
Radiality of the potential gives rotation invariance of its distributional Laplacian.
Let $\nu$ be the image of $\mu_T$ under $z\mapsto|z|^2$.
Averaging the logarithmic kernel over angles yields
\begin{equation}\label{eq:brown-radial-log-kernel}
 f(s)=\frac12\int_{[0,\infty)}\log\max\{s,t\}\dd\nu(t),
 \qquad s>0.
\end{equation}
The scalar identity used here is $(2\pi)^{-1}\int_0^{2\pi}\log|r-te^{i\theta}|\dd\theta =\log\max\{r,t\}$; it follows by expanding the logarithm after factoring out the larger radius, and by an integrable limit when the radii agree.
Compact support bounds the positive part, while the identity bounds the integrated negative part.
Thus the interchange of integrals is valid and the result is finite.

The left and right derivatives of \eqref{eq:brown-radial-log-kernel} are
\[
 f'_-(s)=\frac{\nu([0,s))}{2s},\qquad
 f'_+(s)=\frac{\nu([0,s])}{2s}.
\]
The measure inequality gives $f'_-(s)\le f'_+(s)$, while concavity gives the opposite inequality.
Hence $\nu(\{s\})=0$ for every $s>0$, $f\in C^1((0,\infty))$, and \eqref{eq:brown-radial-cdf} follows.
Also $F(s)/s=2f'(s)$ is nonincreasing.
For $0<s_1<s_2$ this implies
\[
 0\le F(s_2)-F(s_1)
 \le (s_2-s_1)\frac{F(s_1)}{s_1}
 \le\frac{s_2-s_1}{s_1}.
\]
Thus $F$ is locally Lipschitz, with $0\le F'(s)\le F(s)/s$ almost everywhere.
Polar coordinates, $\dd^2z=\tfrac12\dd s\dd\theta$, now give \eqref{eq:brown-radial-density}.
Every possible singular part is concentrated at zero.

The function $F_\varepsilon(z,b)$ in \eqref{eq:brown-regularized-minimum} depends on $z$ only through $s=|z|^2$.
Write $u_\varepsilon(z)=f_\varepsilon(s)$.
Uniqueness of the minimizer makes $b_\varepsilon$ independent of the argument of $z$.
This minimizer is smooth for $s>0$, either from its resolvent expression along $z=\sqrt{s}$ or by differentiation of the stationary equation.
Differentiating the minimum value therefore gives
\[
 f_\varepsilon'(s)=\frac1{2m}
          \Tr(H_{b_\varepsilon(s)}^{-1}b_\varepsilon(s)).
\]
The functions $f_\varepsilon$ are concave and decrease to $f$.
For a fixed positive $s$, bound their derivatives between their left and right secant slopes.
First send $\varepsilon$ to zero, then send the secant increment to zero.
Since $f$ is differentiable, $f_\varepsilon'(s)\to f'(s)$, proving \eqref{eq:brown-radial-optimizer}.

For the spectrum, the critical identities \eqref{eq:boundary-identities} with $a_z=-z1_m$ and $\rho=0$ force $z\ne0$, $Y=-1_m/\overline z$, and
\begin{equation}\label{eq:brown-radial-stein}
 (s\id-\eta)b=1_m,\qquad b>0,\qquad s=|z|^2.
\end{equation}
If $s>r(\eta)$, the positive convergent series $b=\sum_{k\ge0}s^{-k-1}\eta^k(1_m)$ solves this equation.
With the displayed $Y$ it satisfies $\mathcal T_b(Y)=-z1_m$, has $P=H_b^{-1}-Yb^{-1}Y^*=0$, and satisfies \eqref{eq:unregularized-critical}.
Theorem~\ref{thm:main} therefore puts $z$ in the resolvent set.
Conversely, any positive solution of \eqref{eq:brown-radial-stein} makes
\[
 x\longmapsto s^{-1}b^{-1/2}\eta(b^{1/2}xb^{1/2})b^{-1/2}
\]
completely positive with value $1_m-s^{-1}b^{-1}<1_m$ at the identity.
Its norm, and hence its spectral radius, is less than one.
Similarity gives $r(\eta)<s$.
At zero the first critical identity would read $0=Y^{-*}$, which is impossible.
This proves \eqref{eq:brown-spectral-disk}.
Finally Corollary~\ref{cor:brown-capacity} removes the origin atom when capacity is positive.
\end{proof}

The spectral disk in \eqref{eq:brown-spectral-disk} is the full operator spectrum.
The argument above does not show that $\supp\mu_T$ fills that disk.
The estimate on $F(s)/s$ concerns the average mass per unit squared radius; it does not assert that the pointwise density is nonincreasing.
At zero the density need not have a finite bound.

Under the stronger two-sided flatness assumption $c\tr_m(x)1_m\le\eta(x)\le C\tr_m(x)1_m$ for all $x\ge0$, with $C>c>0$, Alt and Kr\"uger prove analyticity and positive upper and lower bounds for the density inside the disk \cite[Proposition~2.9 and Theorem~2.5]{AK20}.
The argument above allows degenerate maps and yields the weaker regularity stated in Propositions~\ref{prop:brown-radial} and~\ref{prop:brown-radial-details}.

\begin{remark}[Comparison with scalar $R$-diagonal elements]
\label{rem:brown-scalar-r-diagonal}
For a scalar $R$-diagonal element, represented in $*$-distribution as $uh$ with $u$ Haar unitary, $h\ge0$ nonconstant and $u,h$ free, Haagerup and Larsen describe the radial distribution through the $S$-transform of the law of $h^2$ \cite[author preprint, Theorem~4.4]{HL00}.
The formula above instead uses the matrix covariance map and its variational minimizers.
The condition $\rho=0$ gives circularity over $M_m(\C)$ in Example~\ref{ex:brown-circular-coefficients}; it does not by itself give scalar $R$-diagonality with respect to $\varphi_m$.
Thus the scalar $S$-transform formula requires an additional hypothesis in the matrix-coefficient setting.
\end{remark}

\subsection{An invertibility criterion and enclosed mass}
The attainment criterion can be expressed using an invertible matrix and a spectral-radius bound.
For $g\in M_m(\C)$ invertible, put
\[
 \Psi_g(x)=g\eta(x)g^*.
\]
% Proof input: Yin_L_F0042, Sections 1--3.
\begin{proposition}\label{prop:brown-stable-resolvent}
A point $z$ lies outside $\spec(T)$ if and only if there is an invertible $g\in M_m(\C)$ such that
\begin{equation}\label{eq:brown-stable-resolvent}
 a_z=g^{-1}+\rho(g),\qquad r(\Psi_g)<1.
\end{equation}
This matrix is unique and equals $g_z=E_m(S_z^{-1})$.
The unique minimizer is
\begin{equation}\label{eq:brown-stable-minimizer}
 b_z=(\id-\Psi_{g_z})^{-1}(g_zg_z^*)
     =\sum_{k\ge0}\Psi_{g_z}^k(g_zg_z^*)>0.
\end{equation}
The strict spectral-radius condition is equivalent to the existence of a positive matrix $b$ with $b-\Psi_g(b)>0$.
\end{proposition}
\begin{proof}
If $z\notin\spec(T)$, Theorem~\ref{thm:main} gives an attained minimum.
The identities \eqref{eq:boundary-identities} then give \eqref{eq:brown-stable-resolvent} with $g=Y^*$ and $b-g\eta(b)g^*=gg^*>0$.
Conjugating by $b^{1/2}$ turns $\Psi_g$ into a completely positive map with value strictly less than $1_m$ at the identity, so its spectral radius is less than one.

Conversely, suppose \eqref{eq:brown-stable-resolvent} holds and define $b$ by the positive series in \eqref{eq:brown-stable-minimizer}.
Then $b=gH_bg^*$.
Put $Y=g^*$.
We have
\[
 H_bYb^{-1}=g^{-1},\qquad
 H_bYb^{-1}+\rho(Y^*)=a_z,
\]
so $Y=\mathcal T_b^{-1}(a_z)$.
Also $H_b^{-1}-Yb^{-1}Y^*=0$ and $Y^*H_bY=b$.
These two identities give \eqref{eq:unregularized-critical}.
Theorem~\ref{thm:main} therefore gives $z\notin\spec(T)$, $g=E_m(S_z^{-1})$ and $b=b_z$.
It also proves uniqueness.

The implication from $b-\Psi_g(b)>0$ to $r(\Psi_g)<1$ has already been proved.
For the reverse implication, $r(\Psi_g)<1$ makes $b=\sum_{k\ge0}\Psi_g^k(1_m)>0$ a solution of $b-\Psi_g(b)=1_m$.
\end{proof}

The equation $a_z=g^{-1}+\rho(g)$ alone need not determine $E_m(S_z^{-1})$.
The additional condition $r(\Psi_g)<1$ selects that matrix; equivalently, one can find $b>0$ with $b-\Psi_g(b)>0$.

\begin{corollary}\label{prop:brown-contour-mass}
Let $\Gamma$ be a positively oriented piecewise smooth simple closed curve in the resolvent set of $T$, and let $\Omega$ be its bounded interior.
Then
\begin{equation}\label{eq:brown-contour-mass}
 \mu_T(\Omega)=-\frac1{2\pi i}\int_\Gamma\tr_m(g_z)\dd z.
\end{equation}
If $\rho=0$, the full spectrum and the exterior potential are
\begin{align}
 \spec(T)&=\spec(a_0)\cup
 \{z\notin\spec(a_0):r(x\mapsto a_z^{-1}\eta(x)a_z^{-*})\ge1\},
                                  \label{eq:brown-affine-circular-spectrum}\\
 u(z)&=\log\Delta_m(a_z)
       =\frac1m\log|\det a_z|,
 \qquad z\notin\spec(T).          \label{eq:brown-exterior-potential}
\end{align}
In this case
\begin{equation}\label{eq:brown-eigenvalue-mass}
 \mu_T(\Omega)=\frac1m
   \#\{\text{eigenvalues of }a_0\text{ in }\Omega,
                        \text{ counted algebraically}\}.
\end{equation}
\end{corollary}
\begin{proof}
By Proposition~\ref{prop:brown-stable-resolvent}, $g_z=E_m(S_z^{-1})$ off the spectrum.
Differentiation of the trace logarithm gives
\[
 2\partial_z u(z)=-\varphi_m(S_z^{-1})=-\tr_m(g_z).
\]
Differentiating \eqref{eq:brown-logarithmic-potential} there instead gives $2\partial_z u(z)=\int(z-w)^{-1}\dd\mu_T(w)$.
The curve has positive distance from $\supp\mu_T$, so integration and the scalar Cauchy formula prove \eqref{eq:brown-contour-mass}.

If $\rho=0$, \eqref{eq:brown-stable-resolvent} forces $g_z=a_z^{-1}$.
The equation has no solution when $a_z$ is singular.
When $a_z$ is invertible, $r(\Psi_{g_z})<1$ holds exactly under the strict inequality complementary to \eqref{eq:brown-affine-circular-spectrum}.
At a resolvent point, $b_z=a_z^{-1}H_{b_z}a_z^{-*}$ and $Z_{b_z}=a_z^{-*}$.
Thus $\det H_{b_z}/\det b_z=|\det a_z|^2$, while the quadratic correction equals $\operatorname{Re}\Tr(a_z^*a_z^{-*})=m$.
Substitution proves \eqref{eq:brown-exterior-potential} in every resolvent component.
Finally $-g_z=(z1_m-a_0)^{-1}$ in \eqref{eq:brown-contour-mass}, and the finite-dimensional argument principle proves \eqref{eq:brown-eigenvalue-mass}.
\end{proof}

In particular, the mass of any region separated by such a contour is a multiple of $1/m$ when $\rho=0$.
The contour hypothesis is part of this statement; it does not give an unconditional formula for the mass of each atom of $\mu_T$.

\begin{remark}[Diagonal variance profiles]
\label{rem:brown-variance-profile}
Assume $\rho=0$, $a_0=\operatorname{diag}(\alpha_1,\ldots,\alpha_m)$ and $\eta(x)=\operatorname{diag}(\sum_j s_{ij}x_{jj})$, with $s_{ij}\ge0$.
Put $\mathsf S=(s_{ij})$ and $D_z=\operatorname{diag}(|\alpha_i-z|^{-2})$ for $z\notin\spec(a_0)$.
The map $x\mapsto a_z^{-1}\eta(x)a_z^{-*}$ vanishes on matrices with zero diagonal and restricts to $D_z\mathsf S$ on diagonal matrices.
Hence its spectral radius equals $r(D_z\mathsf S)$.
In this case \eqref{eq:brown-affine-circular-spectrum} reduces to the variance-profile test for a deformed circular element over a commutative base.
Alt and Kr\"uger use the same threshold $r(\mathsf S D_z)=1$ in \cite[arXiv version, Proposition~5.16]{AK26}; the two products are similar.
Under their block-primitivity assumption A1 and data regularity assumption A2, they also identify the support of the Brown measure and prove boundedness and interior analyticity of the density \cite[arXiv version, Theorem~2.2]{AK26}.
Assumption A2 is automatic for a finite base, whereas A1 remains an additional condition.
The elementary reduction above requires neither condition and concerns the operator spectrum.
\end{remark}

\begin{remark}[Spectral subspaces]
\label{rem:brown-spectral-subspaces}
When the ambient algebra is a $\mathrm{II}_1$-factor, Haagerup and Schultz associate to every Borel set $B$ an invariant spectral subspace whose projection has trace $\mu_T(B)$ \cite[Theorem~1.1]{HS09}.
For a region separated by a resolvent contour, the subspace is the range of the Riesz idempotent $e_\Omega=(2\pi i)^{-1}\int_\Gamma(z1_m-T)^{-1}\dd z$; compare \cite[Introduction, (1.3)]{HS09}.
Thus \eqref{eq:brown-contour-mass} computes the mass of this spectral block.
Under $\rho=0$, the further eigenvalue-count identity follows from the finite-dimensional argument principle.
\end{remark}

\subsection{Density from the variational optimizer}
\label{subsec:brown-density}
To differentiate the minimum value, we first prove that the Hessian at the minimizer is positive definite, using a positive matrix regularizer.
All gradients and Hessians with respect to $b$ are taken on the real Hilbert space of Hermitian matrices with pairing $\langle k,l\rangle=\Tr(kl)$.

% Proof input: Yin_L_F0041, including the matrix-regularizer proof.
\begin{lemma}\label{lem:brown-matrix-regularizer}
Fix an affine pencil $S$ with constant term $a$, and let $\mathcal V_a$ be the function in \eqref{eq:main-action} with $a_0=a$.
For every $D\in M_m(\C)$ with $D>0$,
\begin{equation}\label{eq:brown-matrix-regularizer}
 \tfrac12\varphi_m\log(S^*S+D)
 =\frac1{2m}\min_{b>0}\{\mathcal V_a(b)+\Tr(Db)\}.
\end{equation}
Its unique minimizer is $b_D=E_m((S^*S+D)^{-1})$.
The real Hessian $\mathfrak H_D=D_b^2\mathcal V_a(b_D)$ is positive definite, and its inverse is the completely positive map
\begin{equation}\label{eq:brown-inverse-hessian}
 \mathfrak H_D^{-1}(k)
 =E_m((S^*S+D)^{-1}k(S^*S+D)^{-1}).
\end{equation}
\end{lemma}
\begin{proof}
Apply \eqref{eq:scaled-minimum}, at $\varepsilon=1$, to $SD^{-1/2}$.
Its coefficient maps are
\[
 \eta_D(x)=\eta(D^{-1/2}xD^{-1/2}),\qquad
 \rho_D(x)=\sum_j a_jD^{-1/2}xa_jD^{-1/2}.
\]
Use the trial bijection $b'=D^{1/2}bD^{1/2}$.
If $Z=\mathcal T_b^{-1}(a)$, the corresponding solution for $SD^{-1/2}$ is $Z'=ZD^{1/2}$: its equation is the old equation multiplied on the right by $D^{-1/2}$.
The quadratic correction $\operatorname{Re}\Tr(a^*Z)$ is unchanged, the logarithm of the determinant ratio decreases by $\log\det D$, and the linear term is $\Tr(Db)$.
The same determinant term appears on the operator side, since
\[
 (SD^{-1/2})^*(SD^{-1/2})+1_m
   =D^{-1/2}(S^*S+D)D^{-1/2}.
\]
The trace-log congruence identity follows without commutation: for positive invertible $Q$, differentiate $\varphi_m\log(D^{-t/2}QD^{-t/2})$.
Cyclicity of the trace makes its derivative $-\tr_m\log D$; integrating from $t=0$ to $t=1$ gives the required subtraction of $m^{-1}\log\det D$.
This proves \eqref{eq:brown-matrix-regularizer}.
Transforming the known unique minimizer gives $b_D$ as stated.

Put $Q_D=S^*S+D$.
Stationarity and its derivative in a Hermitian direction $k$ give
\[
 \nabla\mathcal V_a(b_D)+D=0,\qquad
 D_Db_D[k]=-E_m(Q_D^{-1}kQ_D^{-1})=:-\mathcal A_D(k),
 \qquad \mathfrak H_D\mathcal A_D=\id.
\]
The map $\mathcal A_D$ is completely positive and selfadjoint for the trace pairing.
Moreover
\[
 \langle k,\mathcal A_D(k)\rangle
  =(\Tr\otimes\tau)
       ((Q_D^{-1/2}kQ_D^{-1/2})^2)>0
       \qquad(k=k^*\ne0).
\]
Faithfulness and invertibility of $Q_D$ give strictness.
Thus $\mathcal A_D$ is positive definite on the finite-dimensional real space, so $\mathfrak H_D=\mathcal A_D^{-1}>0$ and \eqref{eq:brown-inverse-hessian} follows.
\end{proof}

The same formula at $D=t1_m\downarrow0$ shows that, whenever $S$ is boundedly invertible, the Hessian at its unregularized minimizer is positive definite with inverse $k\mapsto E_m((S^*S)^{-1}k(S^*S)^{-1})$.
This is a statement at the minimizer.
It does not imply global convexity of $\mathcal V_a$.
At positive regularizer one obtains the quadratic-form bounds
\[
 \lambda_{\min}(D)^2\id\le\mathfrak H_D
       \le(\|S\|^2+\|D\|)^2\id.
\]
Indeed the operator $Q_D^{-1/2}kQ_D^{-1/2}$ in the proof has Hilbert--Schmidt norm between $\|Q_D\|^{-1}\|k\|_2$ and $\lambda_{\min}(D)^{-1}\|k\|_2$.
The lower bound can degenerate as $D$ tends to zero.

\begin{proposition}\label{prop:brown-optimizer-density}
For each $\varepsilon>0$, the regularized logarithmic potential satisfies
\begin{equation}\label{eq:brown-envelope-transform}
 \partial_z u_\varepsilon(z)
       =-\frac1{2m}\Tr Z_\varepsilon(z)^*.
\end{equation}
Evaluate all derivatives of $F_\varepsilon$ at $(z,b_\varepsilon(z))$, and set
\[
 \mathfrak H=D_b^2F_\varepsilon>0,\qquad
 h_x=\partial_x\nabla_bF_\varepsilon,
 \qquad h_y=\partial_y\nabla_bF_\varepsilon.
\]
Set $\beta_\varepsilon=(2\pi)^{-1}\nabla_z^2u_\varepsilon$.
Then
\begin{equation}\label{eq:brown-hessian-density}
 \begin{split}
 \beta_\varepsilon(z)=\frac1{4\pi m}\bigl(&
  2\langle1_m,\mathcal T_b^{-1}(1_m)\rangle_{\R}
 +2\langle i1_m,\mathcal T_b^{-1}(i1_m)\rangle_{\R}\\
 &-\langle h_x,\mathfrak H^{-1}h_x\rangle
  -\langle h_y,\mathfrak H^{-1}h_y\rangle\bigr),
 \qquad b=b_\varepsilon(z).
 \end{split}
\end{equation}
It is positive, and $\beta_\varepsilon(z)\dd^2z$ converges weakly to $\mu_T$ as $\varepsilon\downarrow0$.
\end{proposition}
\begin{proof}
For fixed positive $\varepsilon$, the resolvent formulas \eqref{eq:brown-regularized-optimizer} are real analytic in $x,y$.
At the minimizer the derivative with respect to $b$ vanishes, so differentiation of the minimum value uses only the explicit parameter derivative.
The quadratic correction $\operatorname{Re}\Tr(a^*Z_b)$ has first variation $2\langle da,Z_b\rangle_{\R}$.
Since $\partial_xa_z=-1_m$ and $\partial_ya_z=-i1_m$, it follows that
\[
 \partial_xu_\varepsilon=-\frac1m\operatorname{Re}\Tr Z_\varepsilon,
 \qquad
 \partial_yu_\varepsilon=-\frac1m\operatorname{Im}\Tr Z_\varepsilon.
\]
This is \eqref{eq:brown-envelope-transform}.

Lemma~\ref{lem:brown-matrix-regularizer}, with $D=\varepsilon^21_m$, proves $\mathfrak H>0$.
Differentiating stationarity gives $\partial_xb_\varepsilon=-\mathfrak H^{-1}h_x$ and $\partial_yb_\varepsilon=-\mathfrak H^{-1}h_y$.
Hence the real Hessian of the optimized value is the Schur complement
\begin{equation}\label{eq:brown-value-hessian}
 \partial_{jk}u_\varepsilon
  =\frac1{2m}\bigl(F_{jk}
        -\langle h_j,\mathfrak H^{-1}h_k\rangle\bigr),
 \qquad j,k\in\{x,y\}.
\end{equation}
At fixed $b$, $F_\varepsilon(z,b)$ is a real quadratic function of $x,y$, so $F_{xx}=2\langle1_m,\mathcal T_b^{-1}(1_m)\rangle_{\R}$ and $F_{yy}=2\langle i1_m,\mathcal T_b^{-1}(i1_m)\rangle_{\R}$.
Taking the normalized Laplacian proves \eqref{eq:brown-hessian-density}.

For the sign, put $Q=S_z^*S_z+\varepsilon^21_m$ and $P=S_zS_z^*+\varepsilon^21_m$.
Direct trace differentiation and $S_zQ^{-1}S_z^*=1_m-\varepsilon^2P^{-1}$ give
\[
 \partial_{\bar z}\partial_z u_\varepsilon
  =\tfrac12\varphi_m(Q^{-1}-Q^{-1}S_zQ^{-1}S_z^*)
  =\tfrac{\varepsilon^2}{2}\varphi_m(Q^{-1}P^{-1}).
\]
Thus
\begin{equation}\label{eq:positive-density}
 \beta_\varepsilon(z)=\frac{\varepsilon^2}{\pi}
                         \varphi_m(Q^{-1}P^{-1})>0.
\end{equation}
The trace equals that of the positive invertible operator $Q^{-1/2}P^{-1}Q^{-1/2}$; commutation is unnecessary.
The weak convergence follows from \cite[Chapter 11, Section 9, (11.14)--(11.15)]{MS17}, with the normalization in \eqref{eq:brown-definition}.
\end{proof}

\begin{remark}[Free circular addition and uniform density bounds]
\label{rem:brown-circular-addition}
Suppose $x$ is an operator in a finite tracial von Neumann algebra and $c_t$ is a circular element of variance $t>0$, scalar $*$-free from $x$.
Zhong's density formula for $x+c_t$ on its region where the subordination parameter is positive combines a two-resolvent term with the derivative of an implicitly determined parameter \cite[arXiv version, Theorem~4.2]{Zho26}.
This gives a useful comparison with the optimizer differentiation above.
Belinschi, Zhi Yin and Zhong prove, in this setting, that the regularized densities are uniformly bounded by $1/(\pi t)$ and that the limiting Brown measure is absolutely continuous with the same bound \cite[arXiv version, Theorem~7.10 and Lemma~7.11]{BYZ24}.
Such a bound supplies the control needed when the regularizer tends to zero.
The scalar freeness and nonzero circular variance in these results are additional assumptions; the general matrix-coefficient formula \eqref{eq:brown-hessian-density} does not supply them.
\end{remark}

Formula~\eqref{eq:brown-hessian-density} uses the positive minimizer, the inverse of $\mathcal T_b$, and two linear systems with the positive definite coefficient operator $\mathfrak H$ on Hermitian matrices.
It provides another way to compute the same derivatives as the hermitized Dyson equation.
To make the comparison explicit, set
\[
 A_z=\begin{pmatrix}0&a_z\\a_z^*&0\end{pmatrix},\qquad
 G^{-1}=i\varepsilon1_{2m}-A_z-\Theta(G),\qquad
 \operatorname{Im}G<0.
\]
With $E_{21}=\left(\begin{smallmatrix}0&0\\1_m&0\end{smallmatrix}\right)$, differentiation gives the uniquely solvable equation
\begin{equation}\label{eq:linearized-dyson}
 X-G\Theta(X)G=-GE_{21}G,
\end{equation}
and
\begin{equation}\label{eq:potential-derivatives}
 \partial_z u_\varepsilon=\tfrac12\tr_mG_{21},\qquad
 \beta_\varepsilon=\frac1\pi\tr_mX_{21}.
\end{equation}
Indeed $G_{21}=-E_m(Q^{-1}S_z^*)=-Z_\varepsilon^*$ and $X=\partial_{\bar z}G$; invertibility of the linearized equation follows from the weighted contraction after \eqref{eq:imaginary-balance}.
Neither formulation gives a uniform zero-regularizer inverse bound.
Weak convergence permits atoms and curve-supported mass and does not justify pointwise convergence of the densities for arbitrary coefficients.

\begin{remark}[Linear stability of the Dyson equation]\label{rem:dyson-stability}
Writing $\mathcal C_G(X)=GXG$, the operator in \eqref{eq:linearized-dyson} is $\id-\mathcal C_G\Theta$.
Under the convention $M=-G$, it is the stability operator studied in \cite[Section 4.2, Proposition 4.4 and Theorem 2.6]{AEK19}.
Their quantitative estimates use a flatness assumption, and Theorem 2.6 also assumes spatial decay.
Our block covariance does not satisfy the strictly positive lower bound in the flatness assumption: for $P>0$, \eqref{eq:full-covariance} gives
\[
 \Theta\begin{pmatrix}P&0\\0&0\end{pmatrix}
       =\begin{pmatrix}0&0\\0&\eta^*(P)\end{pmatrix}.
\]
Here invertibility at positive $\varepsilon$ follows from the weighted estimate in Section~\ref{sec:dyson}.
At the real parameter used in Section~\ref{sec:main-proof}, the strict inequality \eqref{eq:strict-full-stability} instead gives the invertible linearization needed for analytic continuation and the spectral gap.
\end{remark}

\subsection{Recoveries and examples}
If $a_0,a_j$ are selfadjoint, $T=T^*$ and the smaller matrix Cauchy transform satisfies
\[
 g(w)=E_m((w1_m-T)^{-1}),\qquad
 g(w)^{-1}=w1_m-a_0-\eta(g(w)),\qquad
 \operatorname{Im}g(w)<0\quad(\operatorname{Im}w>0).
\]
Its normalized trace recovers interval masses and atoms by \eqref{eq:stieltjes-interval} and \eqref{eq:stieltjes-atom}; the potential is given by \eqref{eq:selfadjoint-potential}.
For a selfadjoint operator the Brown measure equals its scalar spectral measure on the real line \cite[Chapter 11, p.~270]{MS17}, and faithfulness identifies its support with the spectrum.
The gap criterion in Section~\ref{subsec:gap-tools} therefore applies directly.

For isotropic covariance $\eta(b)=\sigma\tr_m(b)1_m$, Proposition~\ref{prop:minmax-isotropic} reduces the calculation to two real variables.
The density of the spectral measure of $T$ is $q(X^{-1}(x))/(\pi\sigma)$, as stated in \eqref{eq:minmax-isotropic-density}.
This one-dimensional density and the planar densities in Proposition~\ref{prop:brown-radial-details} refer to their respective hypotheses.

For a scalar circular variable of variance $v>0$, put $\eta(b)=vb$ and $\rho=0$.
In \eqref{eq:brown-radial-envelope}, the substitution $t=v+b^{-1}>v$ gives
\[
 f(s)=\frac12\inf_{t>v}\{\log t-1+s/t\}
 =\begin{cases}
   \tfrac12(\log v-1+s/v),&0\le s\le v,\\
   \tfrac12\log s,&s\ge v.
  \end{cases}
\]
For $s\le v$ the infimum is the limit at $t=v$; for $s>v$ it is attained at $t=s$.
Thus $F(s)=\min\{s/v,1\}$ and $\beta(z)=(\pi v)^{-1}1_{\{|z|^2<v\}}$, the normalized uniform law on the disk of radius $\sqrt v$.
If instead $T=\operatorname{diag}(c,0)$ with $c$ standard circular, the normalized potential is the mean of the two diagonal potentials.
It gives half the circular law and half an atom at zero, illustrating the possibility allowed by \eqref{eq:brown-radial-density}.
More generally, invariant block decompositions give weighted sums of the Brown measures of the diagonal blocks; in a $\mathrm{II}_1$-factor this is \cite[Proposition~6.5]{HS09}, with the weights given by the traces of the block projections.

Finally, for the scalar deterministic pencil $T=\lambda$,
\[
 F_\varepsilon(z,b)=-\log b-1+(|\lambda-z|^2+\varepsilon^2)b,
 \qquad b_\varepsilon=(|\lambda-z|^2+\varepsilon^2)^{-1}.
\]
Its Hessian is $b_\varepsilon^{-2}$ and its mixed derivatives are $h_x=2(x-\operatorname{Re}\lambda)$, $h_y=2(y-\operatorname{Im}\lambda)$.
Substitution into \eqref{eq:brown-hessian-density} yields
\[
 \beta_\varepsilon(z)=
 \frac{\varepsilon^2}{\pi(|\lambda-z|^2+\varepsilon^2)^2}.
\]
Polar integration gives total mass one, and the mass outside every fixed neighborhood of $\lambda$ tends to zero.
Hence the limiting measure is $\delta_\lambda$, in agreement with the determinant potential.
This also verifies the signs and both real-direction contributions in the Hessian formula.
